The design of crewed and robotic interplanetary missions requires the simultaneous optimization of
competing objectives across propulsion, trajectory, power, and communications subsystems. This
dissertation presents an integrated framework for mission design that couples high-fidelity trajectory
optimization with multidisciplinary systems analysis to enable rapid trade-space exploration during the
early phases of mission planning.

A variational formulation of the optimal control problem is derived and discretized using a
pseudospectral collocation scheme. The resulting nonlinear program is solved with a sequential
quadratic programming algorithm augmented by a warm-start heuristic based on patched-conic
approximations, yielding an order-of-magnitude reduction in wall-clock solve time compared with
gradient-free methods commonly used in preliminary design.

The framework is validated against published mission data for three reference cases: an Earth-Moon
cargo delivery, a Mars sample-return mission, and an outer solar system flyby sequence. In all cases,
the recovered trajectories match the published delta-V budgets to within \SI{2}{\percent} while
reducing the number of required subsystem evaluations by a factor of four.

Parametric sensitivity analyses reveal that launch vehicle performance has a stronger influence on
total mission cost than either specific impulse or parking orbit altitude for the Mars case, a result
that inverts the conventional design intuition derived from single-objective optimization. The
findings suggest that integrated trajectory-systems co-optimization should be adopted as standard
practice during Phase A mission studies.
